\medskip

Le coyote est un animal sauvage proche du loup, qui vit en Amérique du Nord.

Dans l'état d'Oklahoma, aux États-Unis, 70\,\% des coyotes sont touchés par une maladie appelée ehrlichiose.

Il existe un test aidant à la détection de cette maladie. Lorsque ce test est appliqué à un coyote, son résultat est soit positif, soit négatif, et on sait que :

\setlength\parindent{1cm}
\begin{itemize}
\item[$\bullet~~$] Si le coyote est malade, le test est positif dans 97\,\% des cas.
\item[$\bullet~~$] Si le coyote n'est pas malade, le test est négatif dans 95\,\% des cas.
\end{itemize}
\setlength\parindent{0cm}

\bigskip

\textbf{Partie A}

\medskip

Des vétérinaires capturent un coyote d'Oklahoma au hasard et lui font subir un test pour l'ehrlichiose.

On considère les évènements suivants :

\setlength\parindent{1cm}
\begin{itemize}
\item[$\bullet~~$]$M$: \og le coyote est malade \fg;
\item[$\bullet~~$]$T$: \og le test du coyote est positif \fg.
\end{itemize}
\setlength\parindent{0cm}

On note $\overline{M}$ et $\overline{T}$ respectivement les évènements contraires de $M$ et $T$.

\medskip

\begin{enumerate}
\item Recopier et compléter l'arbre pondéré ci-dessous qui modélise la situation.

\begin{center}
\pstree[treemode=R,nodesepA=0pt,nodesepB=2.5pt,treesep=1cm,levelsep=2.5cm]{\TR{}}
{\pstree{\TR{$M$~}\taput{0,70}}
	{\TR{$T$}\taput{0,97}
	\TR{$\overline{T}$}\tbput{0,03}
	}
\pstree{\TR{$\overline{M}$~}\tbput{0,30}}
	{\TR{$T$}\taput{0,05}
	\TR{$\overline{T}$}\tbput{0,95}
	}
}
\end{center}

\item Déterminer la probabilité que le coyote soit malade et que son test soit positif.

On calcule $P(M \cap T) = P(M) \times P_M(T) = 0,7 \times 0,97 = 0,679$
\item Démontrer que la probabilité de T est égale à $0,694$.

D'après la loi des probabilités totales :

$P(T) = P(M \cap T) + P\left(\overline{M} \cap T\right)$ ; or

$P\left(\overline{M} \cap T\right) = P\left(\overline{M}\right) \times P_{\overline{M}}(T) = 0,3 \times 0,05 = 0,015$, d'où :

$P(T) = 0,679 + 0,015 = 0,694$.
\item On appelle \og valeur prédictive positive du test \fg{} la probabilité que le coyote soit effectivement malade sachant que son test est positif.

Calculer la valeur prédictive positive du test. On arrondira le résultat au millième.

On calcule donc $P_T(M) = \dfrac{P(T \cap M)}{P(T)} = \dfrac{P(M \cap T)}{P(T)} = \dfrac{0,679}{0,694}\approx \np{0,9784}$ soit 0,978 au millième près.
\item  
\begin{enumerate}
\item Par analogie avec la question précédente, proposer une définition de la \og valeur prédictive négative du test \fg{} et calculer cette valeur en arrondissant au millième.

On appelle \og valeur prédictive négative du test \fg{} la probabilité que le coyote ne soit pas malade sachant que son test est négatif.

Elle est égale à $P_{\overline{T}}\left(\overline{M}\right) = \dfrac{P\left(\overline{T} \cap \overline{M} \right)}{P\left(\overline{T}\right)}$.

Or $P\left(\overline{T} \cap \overline{M} \right) = P\left(\overline{T}\right) \times P_{\overline{T}}\left(\overline{M}\right)$ ;

$\bullet~~$$P\left(\overline{T}\right) = 1 - P(T) = 1 - 0,694 = 0,306$ ;

$\bullet~~$$P_{\overline{T}}\left(\overline{M}\right) = \dfrac{P\left(\overline{T} \cap \overline{M}\right)}{P\left(\overline{T}\right)} = \dfrac{0,3 \times 0,95}{0,306} \approx 0,9313$, soit 0,931 au millième près.
\item Comparer les valeurs prédictives positive et négative du test, et interpréter.

Comme $0,978 > 0,931$ cela signifie qu'un résultat positif (erreur d'à peu près 2\,\%) est plus probant qu'un résultat négatif (erreur d'à peu près 7\,\%).
	\end{enumerate}	
\end{enumerate}

\bigskip

\textbf{Partie B}

\medskip

On rappelle que la probabilité qu'un coyote capturé au hasard présente un test positif est de $0,694$.

\medskip

\begin{enumerate}
\item Lorsqu'on capture au hasard cinq coyotes, on assimile ce choix à un tirage avec remise.

On note $X$ la variable aléatoire qui à un échantillon de cinq coyotes capturés au hasard associe le nombre de coyotes dans cet échantillon ayant un test positif.
\begin{enumerate}
\item Quelle est la loi de probabilité suivie par $X$ ? Justifier et préciser ses paramètres.

Le nombre de coyotes est assez important pour que toutes les captures indépendantes sont celles d'animaux dont la probabilité de positivité au test est de 0,694. La variable $X$ suit donc une loi binomiale de paramètres $n = 5$ et $P = 0,694$.
\item Calculer la probabilité que dans un échantillon de cinq coyotes capturés au hasard, un seul ait un test positif. On arrondira le résultat au centième.

On sait que $P(X = 1) = \binom{5}{1} \times 0,694^1 \times (1 - 0,694)^{5 - 1} \approx 0,030$ soit 0,03 au centième près.
\item Un vétérinaire affirme qu'il y a plus d'une chance sur deux qu'au moins quatre coyotes sur cinq aient un test positif : cette affirmation est-elle vraie ? Justifier la réponse.

On vérifie si effectivement $P(X \geqslant 4) > \dfrac{1}{2}$.

Or $P(X \geqslant 4) = P(X = 4) + P(X = 5)$ ;

$P(X = 4) = \binom{5}{4}\times 0,694^4 \times 0,306^1 \approx \np{0,3549}$ et 

$P(X = 5) = \binom{5}{5}\times 0,694^5 \times 0,306^0 \approx \np{0,1609}$

Donc $P(X \geqslant 4) \approx \np{0,3549} + \np{0,3549}$, soit $P(X \geqslant 4) \approx \np{0,5158}$ valeur supérieure à 0,5 : le vétérinaire a raison.
	\end{enumerate}	
\item Pour tester des médicaments, les vétérinaires ont besoin de disposer d'un coyote présentant un test positif. Combien doivent-ils capturer de coyotes pour que la probabilité qu'au moins l'un d'entre eux présente un test positif soit supérieure à $0,99$ ?

Il faut trouver $n \in \N$ tel que $P(Y \geqslant 1) > 0,99$ avec $Y$ variable aléatoire associée au nombre de coyotes ayant un test positif.

Or $P(Y \geqslant 1) = 1 - P(Y = 0) = 1 - \binom{n}{0}\times 0,694^0 \times 0,306^n = 1 - 0,306^n$.
Il faut donc résoudre dans $\N$ l'inéquation : $1 - 0,306^n > 0,99 \iff 0,01 > 0,306^n$

Par croissance de la fonction logarithme népérien : $\ln 0,01 > n \ln 0,306 \iff \dfrac{\ln 0,01}{\ln 0,306} < n$ car $\ln 0,306 < 0$.

Comme $\dfrac{\ln 0,01}{\ln 0,306} \approx 3,9$ : il faut donc capturer au moins 4 coyotes.
\end{enumerate}

\bigskip
